%! Function Composition and Decomposition — Unit 1, Lesson 1.3 — Warm-Up (ANSWER KEY)
% Exactly ONE page, blank and keyed (LESSON_SHAPE.md, one_page). Three prerequisite items
% spiralling 1.1 and 1.2: evaluating at a value and at a negative (and at an EXPRESSION, which
% is the whole mechanic of the BUILDING IT row); a piecewise evaluation; and a function
% operation plus a backwards table read, which is what a handoff lookup needs.
\documentclass[10pt]{article}
\usepackage{precalculus-article}
\usepackage{precalculus-key}
\newcommand{\TallMath}[1]{$\displaystyle #1\rule[-1.4em]{0pt}{3.2em}$}

\begin{document}

\pageheader{Unit 1, Lesson 1.3}{Warm-Up --- Answer Key}

\vspace{0.12in}

\begin{enumerate}[itemsep=12pt]

\item From Lesson 1.1: let $f(x) = 3x - 2$. Whatever you put in the parentheses is what gets
      tripled --- even when it is not a number.

\begin{work}
   f(5) &= 3(5) - 2 = 13 \\
  f(-1) &= 3(-1) - 2 = -5 \\
 f(x^2) &= 3(x^2) - 2 = 3x^2 - 2
\end{work}

\begin{enumerate}[label=(\alph*), itemsep=6pt]
  \item $f(5) = $ \ans{13} \qquad $f(-1) = $ \ans{$-5$}
  \item $f(x^2) = $ \ans{$3x^2 - 2$}
\end{enumerate}

\item From Lesson 1.1: let $v(x) = x^2$ when $x \le 1$, and $v(x) = 4 - x$ when $x > 1$. Decide
      which rule owns the input before you evaluate.

\begin{work}
  v(-2) &= (-2)^2 = 4 \\
   v(1) &= (1)^2 = 1 \\
   v(3) &= 4 - 3 = 1
\end{work}

      $v(-2) = $ \ans{4} \qquad $v(1) = $ \ans{1} \qquad $v(3) = $ \ans{1}

\item From Lesson 1.2: the table gives every value of two functions $f$ and $g$.

\smallskip
\renewcommand{\arraystretch}{1.4}
\begin{center}
\begin{tabular}{|c|c|c|c|c|}
\hline
$x$ & 1 & 2 & 3 & 4 \\
\hline
$f(x)$ & 5 & 2 & 6 & 1 \\
\hline
$g(x)$ & 3 & 4 & 0 & 7 \\
\hline
\end{tabular}
\end{center}

\begin{enumerate}[label=(\alph*), itemsep=6pt]
  \item $(f+g)(2) = $ \ans{6} \qquad $(f-g)(4) = $ \ans{$-6$}
  \item Now read the table the other direction: for which input is $f(x) = 6$? \quad
        $x = $ \ans{3}
\end{enumerate}

\end{enumerate}

\end{document}
